This work presents a demonstrative workflow for automated validation of academic documents produced in LaTeX. The objective is to show, through a short and structurally realistic undergraduate thesis, how to separate content preparation, project compilation, automated checks, and visual review of the final file. The procedure considers objective properties that can be observed in the PDF, including page size, font embedding, conformance with the adopted PDF/A target, and the presence of structural markers. To illustrate result reporting, the study uses a synthetic execution with four checks and a simple pass-rate metric. The presented data are exclusively didactic and do not represent an external experiment. The work also demonstrates practical project organization: metadata concentrated in the main file, separate chapter files, a structured bibliography, and contextual use of a figure, a table, an equation, and source code. The results emphasize that automation can identify repetitive problems earlier and make verification more reproducible, but it does not replace the evaluation of scientific content, methodology, consulted sources, or requirements specific to the applicable course or graduate program.

\keywords{LaTeX; academic works; validation; automation; standardization.}
